\chapter{2013年大纲卷（文）}

\section{单选题}

\begin{problem}
设全集 \(U = \{1, 2, 3, 4, 5\}\)，集合 \(A = \{1, 2\}\)，则 \(\complement_U A =\)（\quad）

\choices
    {\( \{1, 2\} \)}
    {\(\{3,4,5\}\)}
    {\( \{1, 2, 3, 4, 5\} \)}
    {\(\varnothing\)}

\begin{answer}
B
\end{answer}

\begin{solution}
由全集与集合的定义，\(\complement_U A\) 是 \(U\) 中不属于 \(A\) 的元素组成的集合，因此
\(\complement_U A=\{3,4,5\}\). 故选 B.
\end{solution}
\end{problem}

\begin{problem}
已知 \(\alpha\) 是第二象限角，\(\sin \alpha = \frac{5}{13}\)，则 \(\cos \alpha =\)（\quad）

\choices
    {\(-\frac{12}{13}\)}
    {\(-\frac{5}{13}\)}
    {\(\frac{5}{13}\)}
    {\(\frac{12}{13}\)}

\begin{answer}
A
\end{answer}

\begin{solution}
由同角三角函数的基本关系，
\(\cos^2\alpha=1-\sin^2\alpha=1-\left(\frac{5}{13}\right)^2=\frac{144}{169}\)，
故 \(\cos\alpha=\pm\frac{12}{13}\). 因为 \(\alpha\) 是第二象限角，\(\cos\alpha<0\)，所以
\(\cos\alpha=-\frac{12}{13}\). 故选 A.
\end{solution}
\end{problem}

\begin{problem}
已知向量 \(\bs{m} = (\lambda + 1, 1)\)，\(\bs{n} = (\lambda + 2, 2)\)，若 \((\bs{m} + \bs{n}) \perp (\bs{m} - \bs{n})\)，则 \(\lambda =\)（\quad）

\choices
    {\(-4\)}
    {\(-3\)}
    {\(-2\)}
    {\(-1\)}

\begin{answer}
B
\end{answer}

\begin{solution}
由向量垂直的条件，
\[
(\bs{m}+\bs{n})\cdot(\bs{m}-\bs{n})=|\bs{m}|^2-|\bs{n}|^2=0.
\]
于是 \((\lambda+1)^2+1=(\lambda+2)^2+4\)，即
\(\lambda^2+2\lambda+2=\lambda^2+4\lambda+8\)，解得 \(\lambda=-3\). 故选 B.
\end{solution}
\end{problem}

\begin{problem}
不等式 \( |x^2 - 2| < 2 \) 的解集是（\quad）

\choices
    {\((-1, 1)\)}
    {\((-2, 2)\)}
    {\((-1,0)\cup (0,1)\)}
    {\((-2,0)\cup (0,2)\)}

\begin{answer}
D
\end{answer}

\begin{solution}
由绝对值不等式，
\[
-2<x^2-2<2.
\]
两边同时加 \(2\)，得 \(0<x^2<4\). 因而 \(x\ne0\) 且 \(-2<x<2\)，
所以解集为 \((-2,0)\cup(0,2)\). 故选 D.
\end{solution}
\end{problem}

\begin{problem}
\((x + 2)^{8}\) 的展开式中 \(x^{6}\) 的系数是（\quad）

\choices
    {\(28\)}
    {\(56\)}
    {\(112\)}
    {\(224\)}

\begin{answer}
C
\end{answer}

\begin{solution}
由二项式定理，含 \(x^6\) 的项为
\[
\mathrm C_8^2x^6\cdot2^2=112x^6.
\]
因此 \(x^6\) 的系数为 \(112\). 故选 C.
\end{solution}
\end{problem}

\begin{problem}
函数 \( f(x) = \log_2\left(1 + \frac{1}{x}\right)  \) （\(x > 0\)）的反函数 \( f^{-1}(x) = \)（\quad）

\choices
    {\(\frac{1}{2^{x} - 1} \)（\(x > 0\)）}
    {\(\frac{1}{2^{x} - 1} \)（\(x\neq 0\)）}
    {\(2^{x} - 1 \)（\(x \in \R\)）}
    {\(2^{x} - 1 \)（\(x > 0\)）}

\begin{answer}
A
\end{answer}

\begin{solution}
令 \(y=f(x)=\log_2\left(1+\frac{1}{x}\right)\)，其中 \(x>0\). 则
\[
2^y=1+\frac{1}{x},\qquad x=\frac{1}{2^y-1}.
\]
由于 \(x>0\)，有 \(1+\frac{1}{x}>1\)，从而 \(y>0\). 因此反函数的定义域为 \(x>0\)，且
\(f^{-1}(x)=\frac{1}{2^x-1}\). 故选 A.
\end{solution}
\end{problem}

\begin{problem}
已知数列 \( \{a_{n}\} \) 满足 \(3a_{n+1} + a_{n} = 0\)，\(a_{2} = -\frac{4}{3}\)，则 \( \{a_{n}\} \) 的前 10 项和等于（\quad）

\choices
    {\(-6(1 - 3^{-10})\)}
    {\(\frac{1}{9}(1 - 3^{10})\)}
    {\(3(1 - 3^{-10})\)}
    {\(3(1 + 3^{-10})\)}

\begin{answer}
C
\end{answer}

\begin{solution}
由递推关系得 \(a_{n+1}=-\frac{1}{3}a_n\)，所以数列 \(\{a_n\}\) 是以 \(-\frac{1}{3}\) 为公比的等比数列.
由 \(a_2=-\frac{4}{3}\) 得 \(a_1=4\). 因此前 10 项和为
\[
S_{10}=4\frac{1-\left(-\frac{1}{3}\right)^{10}}{1-\left(-\frac{1}{3}\right)}
=3\left(1-3^{-10}\right).
\]
故选 C.
\end{solution}
\end{problem}

\begin{problem}
已知 \(F_{1}(-1,0)\)，\(F_{2}(1,0)\) 是椭圆 \(C\) 的两个焦点，过 \(F_{2}\) 且垂直于 \(x\) 轴的直线交 \(C\) 于 \(A\)，\(B\) 两点，且 \(|AB| = 3\)，则 \(C\) 的方程为（\quad）

\choices
    {\(\frac{x^2}{2} + y^2 = 1\)}
    {\(\frac{x^2}{3} + \frac{y^2}{2} = 1\)}
    {\(\frac{x^2}{4} + \frac{y^2}{3} = 1\)}
    {\(\frac{x^2}{5} + \frac{y^2}{4} = 1\)}

\begin{answer}
C
\end{answer}

\begin{solution}
设椭圆的长半轴、短半轴分别为 \(a\)，\(b\)，焦距参数为 \(c\). 由焦点坐标得 \(c=1\)，
且 \(a^2-b^2=c^2=1\). 直线 \(x=1\) 与椭圆的两个交点为 \(A\)，\(B\)，
由 \(|AB|=3\) 得交点的纵坐标为 \(\pm\frac{3}{2}\). 代入椭圆方程，得
\[
\frac{1}{a^2}+\frac{\left(\frac{3}{2}\right)^2}{b^2}=1.
\]
结合 \(a^2=b^2+1\)，可得
\[
\frac{1}{b^2+1}+\frac{9}{4b^2}=1,
\]
即 \(4b^2+9(b^2+1)=4b^2(b^2+1)\)，从而 \(b^2=3\)，\(a^2=4\).
所以椭圆方程为 \(\frac{x^2}{4}+\frac{y^2}{3}=1\). 故选 C.
\end{solution}
\end{problem}

\begin{problem}
若函数 \( y = \sin (\omega x + \varphi)  \) （\(\omega > 0\)）的部分图像如图所示，则 \( \omega = \)（\quad）

\bitmapfigure[max width=0.58\linewidth]{img/2013/syllabus_liberal/q9_fig1.png}

\choices
    {\(5\)}
    {\(4\)}
    {\(3\)}
    {\(2\)}

\begin{answer}
B
\end{answer}

\begin{solution}
由图象可知，点 \((x_0,y_0)\) 与 \(\left(x_0+\frac{\pi}{4},-y_0\right)\) 的纵坐标互为相反数，
且图象在这两个点之间经历了半个周期，因此
\[
\frac{T}{2}=\frac{\pi}{4},\qquad T=\frac{\pi}{2}.
\]
又 \(T=\frac{2\pi}{\omega}\)，所以 \(\omega=4\). 故选 B.
\end{solution}
\end{problem}

\begin{problem}
已知曲线 \( y = x^4 + ax^2 + 1 \) 在点 \((-1, a + 2)\) 处切线的斜率为 8，则 \( a = \)（\quad）

\choices
    {\(9\)}
    {\(6\)}
    {\(-9\)}
    {\(-6\)}

\begin{answer}
D
\end{answer}

\begin{solution}
对函数求导，得 \(y'=4x^3+2ax\). 在 \(x=-1\) 处切线斜率为
\[
y'\big|_{x=-1}=-4-2a.
\]
由题设 \(-4-2a=8\)，解得 \(a=-6\). 故选 D.
\end{solution}
\end{problem}

\begin{problem}
已知正四棱柱 \(ABCD - A_{1}B_{1}C_{1}D_{1}\) 中，\(AA_{1} = 2AB\)，则 \(CD\) 与平面 \(BDC_{1}\) 所成角的正弦值等于（\quad）

\choices
    {\(\frac{2}{3}\)}
    {\(\frac{\sqrt{3}}{3}\)}
    {\(\frac{\sqrt{2}}{3}\)}
    {\(\frac{1}{3}\)}

\begin{answer}
A
\end{answer}

\begin{solution}
设正方形底面边长为 \(s=AB\)，取空间直角坐标系，使
\(A(0,0,0)\)，\(B(s,0,0)\)，\(C(s,s,0)\)，\(D(0,s,0)\)，\(C_1(s,s,2s)\).
则直线 \(CD\) 的方向向量可取 \(\bs{v}=(1,0,0)\). 平面 \(BDC_1\) 内有两个不共线向量
\(\overrightarrow{DB}=(1,-1,0)s\)，\(\overrightarrow{DC_1}=(1,0,2)s\)，
故其法向量可取
\[
\bs{n}=(1,-1,0)\times(1,0,2)=(-2,-2,1).
\]
设直线 \(CD\) 与平面 \(BDC_1\) 所成角为 \(\theta\)，则
\[
\sin\theta=\frac{|\bs{v}\cdot\bs{n}|}{|\bs{v}|\,|\bs{n}|}
=\frac{2}{1\cdot3}=\frac{2}{3}.
\]
故选 A.
\end{solution}
\end{problem}

\begin{problem}
已知抛物线 \(C: y^{2} = 8x\) 与点 \(M(-2, 2)\)，过 \(C\) 的焦点且斜率为 \(k\) 的直线与 \(C\) 交于 \(A\)，\(B\) 两点，若 \(\overrightarrow{MA} \cdot \overrightarrow{MB} = 0\)，则 \(k =\)（\quad）

\choices
    {\(\frac{1}{2}\)}
    {\( \frac{\sqrt{2}}{2} \)}
    {\(\sqrt{2}\)}
    {\(2\)}

\begin{answer}
D
\end{answer}

\begin{solution}
由 \(y^2=8x=4\cdot2x\)，抛物线的焦点为 \(F(2,0)\). 设直线与抛物线的两个交点分别对应参数
\(t\)，\(u\)，则可写为
\[
A(2t^2,4t),\qquad B(2u^2,4u).
\]
由于 \(A,B,F\) 共线，直线方程为 \(y=k(x-2)\)，代入参数式得
\(4t=2k(t^2-1)\)，\(4u=2k(u^2-1)\). 因而 \(t,u\) 是方程
\(kt^2-2t-k=0\) 的两根，故
\[
t+u=\frac{2}{k},\qquad tu=-1.
\]
又 \(\overrightarrow{MA}=(2t^2+2,4t-2)\)，\(\overrightarrow{MB}=(2u^2+2,4u-2)\)，
由题设垂直得
\[
0=\overrightarrow{MA}\cdot\overrightarrow{MB}
=4\bigl[(t^2+1)(u^2+1)+(2t-1)(2u-1)\bigr].
\]
利用 \(tu=-1\)，并令 \(s=t+u\)，括号内的式子化为
\[
t^2+u^2+2+4tu-2(t+u)+1
=s^2-2tu+2+4tu-2s+1=(s-1)^2.
\]
所以 \(s=1\)，即 \(\frac{2}{k}=1\)，解得 \(k=2\). 故选 D.
\end{solution}
\end{problem}

\section{填空题}

\begin{problem}
设 \( f(x) \) 是以 2 为周期的函数，且当 \( x \in [1,3) \) 时，\( f(x) = x - 2 \)，则 \( f(-1) = \)\fillinblank{}.

\begin{answer}
\(-1\)
\end{answer}

\begin{solution}
因为 \(f(x)\) 的周期为 \(2\)，所以
\(f(-1)=f(-1+2)=f(1)\). 由题设在区间 \([1,3)\) 上有 \(f(x)=x-2\)，
故 \(f(1)=1-2=-1\)，从而 \(f(-1)=-1\).
\end{solution}
\end{problem}

\begin{problem}
从进入决赛的6名选手中决出1名一等奖，2名二等奖，3名三等奖，则可能的决赛结果共有\fillinblank{} 种.

\begin{answer}
60
\end{answer}

\begin{solution}
先从 6 名选手中选出 1 名一等奖获得者，再从剩下的 5 名中选出 2 名二等奖获得者，
其余 3 人获得三等奖. 因此可能的决赛结果共有
\[
\mathrm C_6^1\mathrm C_5^2=6\times10=60
\]
种.
\end{solution}
\end{problem}

\begin{problem}
若 \(x\)，\(y\) 满足的约束条件 \(\begin{cases}
x\geqslant 0,\\
x + 3y\geqslant 4,\\
3x + y\leqslant 4.
\end{cases}\) 则 \(z = -x + y\) 的最小值为\fillinblank{}.

\begin{answer}
0
\end{answer}

\begin{solution}
约束条件确定的可行域顶点由边界直线两两相交得到. 当 \(x=0\) 时，
由 \(x+3y=4\) 和 \(3x+y=4\) 分别得到顶点 \(\left(0,\frac{4}{3}\right)\) 和 \((0,4)\)；
两条斜边界直线的交点满足
\[
x+3y=4,\qquad 3x+y=4,
\]
解得 \((x,y)=(1,1)\). 目标函数 \(z=-x+y\) 在三个顶点处的值分别为
\(\frac{4}{3}\)，\(4\)，\(0\). 因而 \(z\) 的最小值为 \(0\).
\end{solution}
\end{problem}

\begin{problem}
已知圆 \(O\) 和圆 \(K\) 是球 \(O\) 的大圆和小圆，其公共弦长等于球 \(O\) 的半径，\(OK = \frac{3}{2}\)，且圆 \(O\) 与圆 \(K\) 所在的平面所成的一个二面角为 \(60^{\circ}\)，则球 \(O\) 的表面积等于\fillinblank{}.

\begin{answer}
\(16\pi\)
\end{answer}

\begin{solution}
设球 \(O\) 的半径为 \(R\)，两圆公共弦的中点为 \(H\). 公共弦长为 \(R\)，
所以半弦长为 \(\frac{R}{2}\). 因为圆 \(O\) 是大圆，其半径为 \(R\)，在大圆所在平面内有
\[
OH=\sqrt{R^2-\left(\frac{R}{2}\right)^2}=\frac{\sqrt{3}}{2}R.
\]
圆 \(K\) 所在平面到球心的距离为 \(OK=\frac{3}{2}\)，故圆 \(K\) 的半径满足
\(r_K^2=R^2-\frac{9}{4}\). 同理，在圆 \(K\) 所在平面内有
\[
KH=\sqrt{r_K^2-\left(\frac{R}{2}\right)^2}
=\frac{\sqrt{3}}{2}\sqrt{R^2-3}.
\]
两平面所成二面角为 \(60^{\circ}\)，而 \(OH\)，\(KH\) 分别垂直于公共弦，
故 \(\angle OHK=60^{\circ}\). 在三角形 \(OHK\) 中由余弦定理，
\[
\left(\frac{3}{2}\right)^2
=\left(\frac{\sqrt{3}}{2}R\right)^2
+\left(\frac{\sqrt{3}}{2}\sqrt{R^2-3}\right)^2
-2\cdot\frac{\sqrt{3}}{2}R\cdot\frac{\sqrt{3}}{2}\sqrt{R^2-3}\cdot\frac{1}{2}.
\]
化简得 \(R\sqrt{R^2-3}=2(R^2-3)\). 由于 \(R^2>3\)，可约去 \(\sqrt{R^2-3}\)，
得到 \(R=2\sqrt{R^2-3}\)，从而 \(R^2=4\). 因此球的表面积为
\(4\pi R^2=16\pi\).
\end{solution}
\end{problem}

\section{解答题}

\begin{problem}
等差数列 \(\{a_{n}\}\) 中，\(a_{7} = 4\)，\(a_{19} = 2a_{9}\).

\begin{enumerate}
\item 求 \(\{a_{n}\}\) 的通项公式；
\item 设 \(b_n = \frac{1}{na_n}\)，求数列 \(\{b_n\}\) 的前 \(n\) 项和 \(S_n\).
\end{enumerate}
\end{problem}

\begin{solution}
\begin{enumerate}
\item 设等差数列 \(\{a_n\}\) 的首项为 \(a_1\)，公差为 \(d\)，则
\(a_7=a_1+6d=4\)，且 \(a_{19}=a_1+18d\)，\(a_9=a_1+8d\).
由 \(a_{19}=2a_9\) 得 \(a_1+18d=2(a_1+8d)\)，所以 \(a_1=2d\).
代入 \(a_1+6d=4\)，得 \(8d=4\)，即 \(d=\frac{1}{2}\)，\(a_1=1\).
因此
\(a_n=a_1+(n-1)d=1+\frac{n-1}{2}=\frac{n+1}{2}\).

\item 由上问知 \(a_n=\frac{n+1}{2}\)，所以
\[
b_n=\frac{1}{na_n}=\frac{2}{n(n+1)}=2\left(\frac{1}{n}-\frac{1}{n+1}\right).
\]
于是
\[
S_n=2\left(1-\frac{1}{2}+\frac{1}{2}-\frac{1}{3}+\cdots+\frac{1}{n}-\frac{1}{n+1}\right)=2\left(1-\frac{1}{n+1}\right)=\frac{2n}{n+1}.
\]
\end{enumerate}
\end{solution}

\begin{problem}
设 \(\triangle ABC\) 的内角 \(A\)，\(B\)，\(C\) 的对边分别为 \(a\)，\(b\)，\(c\)，\((a + b + c)(a - b + c) = ac\).

\begin{enumerate}
\item 求 \(B\)；
\item 若 \(\sin A \sin C = \frac{\sqrt{3} - 1}{4}\)，求 \(C\).
\end{enumerate}
\end{problem}

\begin{solution}
\begin{enumerate}
\item 由题设，\((a+b+c)(a-b+c)=(a+c)^2-b^2\).
根据余弦定理，\(b^2=a^2+c^2-2ac\cos B\)，所以
\[
(a+c)^2-b^2=a^2+2ac+c^2-a^2-c^2+2ac\cos B=2ac(1+\cos B).
\]
由题设等式得 \(2ac(1+\cos B)=ac\).
由于 \(a>0\)，\(c>0\)，故 \(2(1+\cos B)=1\)，即 \(\cos B=-\frac{1}{2}\).
又 \(0<B<\pi\)，所以 \(B=\frac{2\pi}{3}\).

\item 由三角形内角和得 \(A+C=\frac{\pi}{3}\). 利用积化和差公式，
\[
\sin A\sin C=\frac{\cos(A-C)-\cos(A+C)}{2}=\frac{\cos(A-C)-\frac{1}{2}}{2}=\frac{\sqrt{3}-1}{4},
\]
所以 \(\cos(A-C)=\frac{\sqrt{3}}{2}\).
因为 \(-\frac{\pi}{3}<A-C<\frac{\pi}{3}\)，故 \(A-C=\pm\frac{\pi}{6}\).
当 \(A-C=\frac{\pi}{6}\) 时，\(C=\frac{\pi}{12}\)；当 \(A-C=-\frac{\pi}{6}\) 时，\(C=\frac{\pi}{4}\).
因此 \(C=\frac{\pi}{12}\) 或 \(C=\frac{\pi}{4}\).
\end{enumerate}
\end{solution}

\begin{problem}
如图，四棱锥 \(P - ABCD\) 中，\(\angle ABC = \angle BAD = 90^{\circ}\)，\(BC = 2AD\)，\(\triangle PAB\) 与 \(\triangle PAD\) 都是边长为 2 的等边三角形.

\begin{enumerate}
\item 证明：\(PB \perp CD\)；
\item 求点 \(A\) 到平面 \(PCD\) 的距离.
\end{enumerate}

\FigureLayoutDeclare{img/2013/syllabus_liberal/q19_fig1.png}{7.0cm}{16bp 23bp 102bp 67bp}
\bitmapfigure[max width=0.42\linewidth]{img/2013/syllabus_liberal/q19_fig1.png}
\end{problem}

\begin{solution}
\begin{enumerate}
\item 由 \(\triangle PAB\) 和 \(\triangle PAD\) 都是边长为 2 的等边三角形，得 \(AB=AD=2\)，从而 \(BC=4\).
建立空间直角坐标系，使底面 \(ABCD\) 在 \(z=0\) 平面内，并取
\(A(0,0,0)\)，\(B(2,0,0)\)，\(D(0,2,0)\). 按四边形顶点顺序可取 \(C(2,4,0)\).
设 \(P(x,y,z)\)，由 \(PA=PB\) 得 \(x=1\)，由 \(PA=PD\) 得 \(y=1\).
又 \(PA=2\)，所以 \(1^2+1^2+z^2=4\). 取点 \(P\) 在底面上方，得 \(P(1,1,\sqrt{2})\).
于是 \(\overrightarrow{BP}=(-1,1,\sqrt{2})\)，\(\overrightarrow{DC}=(2,2,0)\)，从而
\(\overrightarrow{BP}\cdot\overrightarrow{DC}=-2+2=0\). 因此 \(PB\perp CD\).

\item 由 \(\overrightarrow{DC}=(2,2,0)\)，\(\overrightarrow{DP}=(1,-1,\sqrt{2})\)，可取平面 \(PCD\) 的一个法向量为
\[
\overrightarrow{DC}\times\overrightarrow{DP}=(2\sqrt{2},-2\sqrt{2},-4),
\]
即取 \(\bs{n}=(\sqrt{2},-\sqrt{2},-2)\).
因此平面 \(PCD\) 的方程为 \(\sqrt{2}x-\sqrt{2}(y-2)-2z=0\).
点 \(A(0,0,0)\) 到该平面的距离为
\[
\frac{|\sqrt{2}\cdot0-\sqrt{2}\cdot0-2\cdot0+2\sqrt{2}|}{\sqrt{(\sqrt{2})^2+(-\sqrt{2})^2+(-2)^2}}=\frac{2\sqrt{2}}{2\sqrt{2}}=1.
\]
所以点 \(A\) 到平面 \(PCD\) 的距离为 \(1\).
\end{enumerate}
\end{solution}

\begin{problem}
甲，乙，丙三人进行羽毛球练习赛，其中两人比赛，另一人当裁判. 每局比赛结束时，负的一方在下一局当裁判. 设各局中双方获胜的概率均为 \(\frac{1}{2}\)，各局比赛的结果都相互独立，第1局甲当裁判.

\begin{enumerate}
\item 求第 4 局甲当裁判的概率；
\item 求前 4 局中乙恰好当 1 次裁判的概率.
\end{enumerate}
\end{problem}

\begin{solution}
\begin{enumerate}
\item 第 1 局甲当裁判，所以第 2 局的裁判只能是乙或丙，且概率各为 \(\frac{1}{2}\).
若第 2 局乙当裁判，要使第 4 局甲当裁判，必须依次出现第 3 局丙当裁判、第 4 局甲当裁判，
该路径的概率为 \(\frac{1}{2}\cdot\frac{1}{2}\cdot\frac{1}{2}=\frac{1}{8}\).
若第 2 局丙当裁判，同理必须依次出现第 3 局乙当裁判、第 4 局甲当裁判，概率也是 \(\frac{1}{8}\).
因此第 4 局甲当裁判的概率为 \(\frac{1}{8}+\frac{1}{8}=\frac{1}{4}\).

\item 第 1 局甲当裁判，乙在前 4 局中恰好当 1 次裁判的可能路径分两类.
若第 2 局乙当裁判，则第 3、4 局裁判必须依次为甲、丙或丙、甲，共有 2 条路径；
若第 2 局丙当裁判，则第 3、4 局裁判可以依次为甲、乙，乙、甲或乙、丙，共有 3 条路径.
因此共有 5 条路径. 每条路径需要连续 3 次比赛结果确定，概率均为
\(\left(\frac{1}{2}\right)^3=\frac{1}{8}\)，故所求概率为
\(5\cdot\frac{1}{8}=\frac{5}{8}\).
\end{enumerate}
\end{solution}

\begin{problem}
已知函数 \( f(x) = x^3 + 3ax^2 + 3x + 1 \).

\begin{enumerate}
\item 当 \(a = -\sqrt{2}\) 时，讨论 \(f(x)\) 的单调性；
\item 若 \(x \in [2, +\infty)\) 时，\(f(x) \geqslant 0\)，求 \(a\) 的取值范围.
\end{enumerate}
\end{problem}

\begin{solution}
\begin{enumerate}
\item 当 \(a=-\sqrt{2}\) 时，
\[
f'(x)=3x^2-6\sqrt{2}x+3
=3(x^2-2\sqrt{2}x+1)
=3\bigl(x-(\sqrt{2}-1)\bigr)\bigl(x-(\sqrt{2}+1)\bigr).
\]
因为 \(\sqrt{2}-1<\sqrt{2}+1\)，所以当 \(x<\sqrt{2}-1\) 或 \(x>\sqrt{2}+1\) 时，\(f'(x)>0\)；
当 \(\sqrt{2}-1<x<\sqrt{2}+1\) 时，\(f'(x)<0\).
因此 \(f(x)\) 在 \((-\infty,\sqrt{2}-1)\) 和 \((\sqrt{2}+1,+\infty)\) 上单调递增，
在 \((\sqrt{2}-1,\sqrt{2}+1)\) 上单调递减.

\item 因为 \(2\in[2,+\infty)\)，由 \(f(2)\geqslant0\) 得 \(15+12a\geqslant0\)，所以 \(a\geqslant-\frac{5}{4}\).
反之，当 \(a\geqslant-\frac{5}{4}\) 且 \(x\geqslant2\) 时，
\[
f'(x)=3(x^2+2ax+1)\geqslant3\left(x^2-\frac{5}{2}x+1\right)=3(x-2)\left(x-\frac{1}{2}\right)\geqslant0.
\]
所以 \(f(x)\) 在 \([2,+\infty)\) 上单调不减，从而 \(f(x)\geqslant f(2)=15+12a\geqslant0\).
因此 \(a\) 的取值范围是 \(\left[-\frac{5}{4},+\infty\right)\).
\end{enumerate}
\end{solution}

\section{选考题：解答题}

\begin{problem}
已知双曲线 \(C: \frac{x^2}{a^2} - \frac{y^2}{b^2} = 1 \) （\(a > 0, b > 0\)）的左，右焦点分别为 \(F_1\)，\(F_2\)，离心率为 3，直线 \(y = 2\) 与 \(C\) 的两个交点间的距离为 \(\sqrt{6}\).

\begin{enumerate}
\item 求 \(a\)，\(b\)；
\item 设过 \(F_{2}\) 的直线 \(l\) 与 \(C\) 的左，右两支分别交于 \(A\)，\(B\) 两点，且 \(|AF_{1}| = |BF_{1}|\)，证明：\(|AF_{2}|\)，\(|AB|\)，\(|BF_{2}|\) 成等比数列.
\end{enumerate}
\end{problem}

\begin{solution}
\begin{enumerate}
\item 双曲线的离心率为 \(3\)，所以 \(\frac{c}{a}=3\)，即 \(c=3a\).
又双曲线满足 \(c^2=a^2+b^2\)，故 \(b^2=8a^2\).
直线 \(y=2\) 与双曲线的交点横坐标互为相反数，设右侧交点横坐标为 \(x_0>0\).
由两交点间距离为 \(\sqrt{6}\) 得 \(2x_0=\sqrt{6}\)，即 \(x_0^2=\frac{3}{2}\).
将 \(y=2\) 代入双曲线方程，得 \(\frac{x_0^2}{a^2}-\frac{4}{b^2}=1\).
代入 \(x_0^2=\frac{3}{2}\) 及 \(b^2=8a^2\)，有 \(\frac{3}{2a^2}-\frac{1}{2a^2}=1\)，所以 \(a^2=1\).
由于 \(a>0\)，得 \(a=1\)，进而 \(b^2=8\)，\(b=2\sqrt{2}\).

\item 由上问，双曲线方程为 \(x^2-\frac{y^2}{8}=1\)，且 \(F_1(-3,0)\)，\(F_2(3,0)\).
设 \(A(x_1,y_1)\)，\(B(x_2,y_2)\)，其中 \(x_1<-1\)，\(x_2>1\).
对于双曲线上的任意点 \(P(x,y)\)，有
\[
PF_1^2=(x+3)^2+y^2=(x+3)^2+8(x^2-1)=(3x+1)^2.
\]
因此 \(AF_1=-(3x_1+1)\)，\(BF_1=3x_2+1\).
由 \(AF_1=BF_1\) 得 \(-3x_1-1=3x_2+1\)，即 \(x_1+x_2=-\frac{2}{3}\).

设直线 \(l\) 的方程为 \(y=k(x-3)\).
竖直直线 \(x=3\) 只与双曲线的右支相交，不符合题设，故可作此设.
令 \(m=\frac{k^2}{8}\)，将直线方程代入双曲线方程，得
\[
x^2-m(x-3)^2=1,
\]
即 \((1-m)x^2+6mx-(9m+1)=0\).
其两根为 \(x_1,x_2\)，所以 \(x_1+x_2=-\frac{6m}{1-m}=\frac{6m}{m-1}\).
结合 \(x_1+x_2=-\frac{2}{3}\)，得 \(m=\frac{1}{10}\)，从而 \(k^2=\frac{4}{5}\).
此时根满足 \(9x^2+6x-19=0\)，故
\[
x_1=\frac{-1-2\sqrt{5}}{3},\qquad x_2=\frac{-1+2\sqrt{5}}{3}.
\]
由 \(x_1<-1<x_2<3\)，在线段所在直线上点的顺序为 \(A\)，\(B\)，\(F_2\).
又
\[
AF_2=1-3x_1=2+2\sqrt{5},\qquad BF_2=3x_2-1=2\sqrt{5}-2.
\]
于是 \(AB=AF_2-BF_2=4\)，并且
\[
AB^2=16=(2+2\sqrt{5})(2\sqrt{5}-2)=AF_2\cdot BF_2.
\]
所以 \(|AF_2|\)，\(|AB|\)，\(|BF_2|\) 成等比数列.
\end{enumerate}
\end{solution}
